\documentclass[11pt]{article}
\usepackage{../common/pbstyle}
\input{../common/sheethead}
\renewcommand{\course}{Ensemble Methods -- TD PAC-Bayes}
%% \soltrue is set by the solution wrapper (td_pacbayes_solutions.tex)
\providecommand{\ifsol}{\iffalse}
\usepackage{comment}
\ifsol
\newtcolorbox{solution}{enhanced,breakable,colback=pbteal!4,colframe=pbteal,boxrule=0.4pt,sharp corners,title={Solution},fonttitle=\bfseries\small,colbacktitle=pbteal!15,coltitle=black}
\else
\excludecomment{solution}
\fi
\newcounter{exo}
\usepackage{needspace}
\newenvironment{exo}[1]{\Needspace{8\baselineskip}\refstepcounter{exo}\medskip\noindent{\large\bfseries\color{pbblue}Exercise \theexo\ -- #1}\par\smallskip}{\par}

\begin{document}
\sloppy
\sheethead{TD -- PAC-Bayesian Learning\ifsol\ -- Solutions\fi}{This tutorial accompanies Sections~2--7 of the lecture notes
\emph{PAC-Bayesian Learning: a Theoretical Framework for Ensemble Methods}. The exercises are
of increasing difficulty; those marked $(\star)$ are at Master level. Notation: $S\sim\D^m$,
voters $h:\X\to\{-1,+1\}$, posterior $Q$, prior $P$, Gibbs classifier $\GQ$, majority vote $\BQ$,
$W_Q(x,y)=\E_{h\sim Q}\ind{h(x)\neq y}$, joint error $e_Q=\E W_Q^2$, disagreement $d_Q$.}

%% =====================================================================
\begin{exo}{The Kullback--Leibler divergence}\label{exo:kl}
Let $\Hcal=\{h_1,\dots,h_n\}$ be finite and $P,Q$ two distributions on $\Hcal$ with $P_i>0$.
\begin{enumerate}
\item Using Jensen's inequality, show that $\KL(Q\|P)\geq0$, with equality iff $Q=P$.
\item Assume $P$ uniform. Show that $\KL(Q\|P)=\ln n-H(Q)$ where $H(Q)=-\sum_iQ_i\ln Q_i$. What are the
      minimal and maximal values of $\KL(Q\|P)$? Compute it when $Q$ is uniform over $k\leq n$ voters.
\item Show that $\KL(Q\otimes Q\,\|\,P\otimes P)=2\KL(Q\|P)$.
\item Show that $\KL\big(\mathcal N(\mu_1,\sigma^2I_d)\,\|\,\mathcal N(\mu_0,\sigma^2I_d)\big)=\frac{\|\mu_1-\mu_0\|^2}{2\sigma^2}$.
\item $(\star)$ Binary kl. Fix $q\in[0,1]$. Show that $p\mapsto\kl(q\|p)$ is convex on $(0,1)$, minimal at $p=q$, and
      prove Pinsker's inequality $\kl(q\|p)\geq2(q-p)^2$ by studying $g(p)=\kl(q\|p)-2(q-p)^2$.
\end{enumerate}
\end{exo}

\begin{solution}
\begin{enumerate}
\item \emph{Idea: $-\KL$ is the expectation of a logarithm, and $\ln$ is concave.} Let $I=\{i:Q_i>0\}$ be the
      support of $Q$. By definition, $-\KL(Q\|P)=\sum_{i\in I}Q_i\ln\frac{P_i}{Q_i}$. Since the $Q_i$, $i\in I$, are
      weights summing to one and $\ln$ is concave, Jensen's inequality gives
      \[-\KL(Q\|P)\leq\ln\Big(\sum_{i\in I}Q_i\frac{P_i}{Q_i}\Big)=\ln\Big(\sum_{i\in I}P_i\Big)\leq\ln1=0.\]
      \emph{Equality case.} Jensen's inequality for a strictly concave function is an equality iff the variable is
      constant: $P_i/Q_i=c$ for all $i\in I$. The second inequality is an equality iff $\sum_{i\in I}P_i=1$, i.e.\
      $I=\Hcal$ (since all $P_i>0$). Then $c=\sum_iP_i/\sum_iQ_i=1$ and $Q=P$. Conversely $\KL(P\|P)=0$.
\item With $P_i=\frac1n$: $\KL(Q\|P)=\sum_iQ_i\ln(nQ_i)=\ln n\sum_iQ_i+\sum_iQ_i\ln Q_i=\ln n-H(Q)$.
      The entropy satisfies $H(Q)\geq0$ (each term $-Q_i\ln Q_i\geq0$), with equality for a Dirac, and
      $H(Q)\leq\ln n$ (this is question~1 applied to $\KL(Q\|\text{uniform})\geq0$), with equality for the uniform
      distribution. Hence $0\leq\KL(Q\|P)\leq\ln n$: the KL is $0$ when the posterior is the prior and $\ln n$ when it
      selects a single voter (the Occam price). For $Q$ uniform on $k$ voters, $H(Q)=k\cdot\frac1k\ln k=\ln k$ and
      $\KL=\ln\frac nk$: selecting a tenth of the voters costs $\ln10\approx2.3$ nats, whatever $n$.
\item Write $(Q\otimes Q)_{ij}=Q_iQ_j$. Then $\ln\frac{Q_iQ_j}{P_iP_j}=\ln\frac{Q_i}{P_i}+\ln\frac{Q_j}{P_j}$, so
      \begin{align*}\KL(Q^2\|P^2)&=\sum_{i,j}Q_iQ_j\ln\frac{Q_i}{P_i}+\sum_{i,j}Q_iQ_j\ln\frac{Q_j}{P_j}\\
      &=\sum_iQ_i\ln\frac{Q_i}{P_i}+\sum_jQ_j\ln\frac{Q_j}{P_j}=2\KL(Q\|P),\end{align*}
      using $\sum_jQ_j=1$ in the first sum and $\sum_iQ_i=1$ in the second. This is why second-order bounds
      (tandem, disagreement) pay a doubled KL.
\item The two densities have the same normalizing constant, so the log-ratio is
      $\ln\frac{q(\theta)}{p(\theta)}=\frac{\|\theta-\mu_0\|^2-\|\theta-\mu_1\|^2}{2\sigma^2}$. Write
      $\theta=\mu_1+\sigma\varepsilon$ with $\varepsilon\sim\mathcal N(0,I_d)$ under $Q$. Then
      $\|\theta-\mu_1\|^2=\sigma^2\|\varepsilon\|^2$ and
      $\|\theta-\mu_0\|^2=\|\mu_1-\mu_0\|^2+2\sigma\langle\varepsilon,\mu_1-\mu_0\rangle+\sigma^2\|\varepsilon\|^2$. The
      terms in $\|\varepsilon\|^2$ cancel, and $\E\langle\varepsilon,\mu_1-\mu_0\rangle=0$, so
      $\KL=\E_Q\ln\frac qp=\frac{\|\mu_1-\mu_0\|^2}{2\sigma^2}$. Moving the centre of a Gaussian posterior costs the
      squared distance travelled, measured in units of the variance.
\item For $p\in(0,1)$, $\partial_p\kl(q\|p)=-\frac qp+\frac{1-q}{1-p}=\frac{-q(1-p)+p(1-q)}{p(1-p)}=\frac{p-q}{p(1-p)}$,
      which is negative for $p<q$ and positive for $p>q$: the function decreases then increases, and is minimal at
      $p=q$ where it vanishes. Moreover $\partial_p^2\kl(q\|p)=\frac q{p^2}+\frac{1-q}{(1-p)^2}>0$ (unless $q\in\{0,1\}$,
      where it is still $\geq0$): the function is convex.
      \emph{Pinsker.} Let $g(p)=\kl(q\|p)-2(q-p)^2$. Then $g(q)=0$ and
      \[g'(p)=\frac{p-q}{p(1-p)}-4(p-q)=(p-q)\Big(\frac{1}{p(1-p)}-4\Big).\]
      Since $p(1-p)\leq\frac14$, the bracket is non-negative, so $g'$ has the sign of $p-q$: $g$ decreases on $(0,q)$,
      increases on $(q,1)$, and its minimum is $g(q)=0$. Hence $\kl(q\|p)\geq2(q-p)^2$. The inequality is tight
      when $p(1-p)\approx\frac14$, i.e.\ around $p=\frac12$, which is why the kl and Hoeffding bounds coincide there.
\end{enumerate}
\end{solution}

%% =====================================================================
\begin{exo}{Change of measure and Gibbs posterior}\label{exo:dv}
Let $\phi:\Hcal\to\R$ and $P$ a distribution on a finite set $\Hcal$. Define $Q_\phi(h)=\frac{P(h)e^{\phi(h)}}{\E_{P}e^{\phi}}$.
\begin{enumerate}
\item Show that for any $Q$, $\KL(Q\|Q_\phi)=\KL(Q\|P)-\E_Q\phi+\ln\E_Pe^\phi$.
\item Deduce the change of measure inequality $\E_Q\phi\leq\KL(Q\|P)+\ln\E_Pe^\phi$ and the variational formula
      $\ln\E_Pe^{\phi}=\max_Q\{\E_Q\phi-\KL(Q\|P)\}$.
\item Let $\lambda>0$. Show that the minimizer of $Q\mapsto\RS(\GQ)+\frac{1}{\lambda m}\KL(Q\|P)$ is the Gibbs posterior
      $Q_\lambda(h)\propto P(h)e^{-\lambda m\RS(h)}$, and give the minimal value.
\item Describe $Q_\lambda$ when $\lambda\to0$ and $\lambda\to+\infty$. Interpret $\lambda$.
\item Recover $Q_\lambda$ directly with a Lagrange multiplier for the constraint $\sum_hQ(h)=1$.
\end{enumerate}
\end{exo}

\begin{solution}
\begin{enumerate}
\item By definition of $Q_\phi$, $\ln\frac{Q(h)}{Q_\phi(h)}=\ln\frac{Q(h)}{P(h)}-\phi(h)+\ln\E_Pe^\phi$ (the last term does
      not depend on $h$). Taking the expectation under $Q$:
      $\KL(Q\|Q_\phi)=\KL(Q\|P)-\E_Q\phi+\ln\E_Pe^\phi$.
\item By Exercise~\ref{exo:kl}, $\KL(Q\|Q_\phi)\geq0$, hence $\E_Q\phi\leq\KL(Q\|P)+\ln\E_Pe^\phi$ for every $Q$: this is
      the change of measure inequality. It is an equality iff $\KL(Q\|Q_\phi)=0$, i.e.\ iff $Q=Q_\phi$. Therefore
      $\E_Q\phi-\KL(Q\|P)\leq\ln\E_Pe^\phi$ for all $Q$, with equality at $Q_\phi$: the maximum over $Q$ is reached
      and equals $\ln\E_Pe^\phi$. \emph{Interpretation:} the left-hand side of the change of measure concerns a
      posterior that may depend on the data, the right-hand side only involves the prior; the price of the
      transfer is exactly the KL.
\item Apply question~2 to $\phi=-\lambda m\RS$ (a function of $h$ once $S$ is fixed):
      $-\lambda m\,\E_Q\RS-\KL(Q\|P)\leq\ln\E_Pe^{-\lambda m\RS}$, i.e., dividing by $-\lambda m<0$,
      \[\RS(\GQ)+\frac{\KL(Q\|P)}{\lambda m}\geq-\frac1{\lambda m}\ln\E_{h\sim P}e^{-\lambda m\RS(h)},\]
      with equality iff $Q=Q_\phi$, i.e.\ $Q(h)\propto P(h)e^{-\lambda m\RS(h)}$. The minimal value is the right-hand side,
      a ``soft minimum'' of the empirical risks (it tends to $\min_h\RS(h)$ when $\lambda\to\infty$ on a finite class).
\item When $\lambda\to0$, $e^{-\lambda m\RS(h)}\to1$ for every $h$ and $Q_\lambda\to P$: the data are ignored.
      When $\lambda\to\infty$, divide numerator and denominator by $e^{-\lambda m r^\star}$ with $r^\star=\min_h\RS(h)$:
      the weight of every $h$ with $\RS(h)>r^\star$ tends to $0$, and $Q_\lambda$ converges to $P$ restricted to the
      empirical risk minimizers and renormalized. $\lambda$ is an inverse temperature: it sets how much a
      difference of one training error changes the relative weights (by a factor $e^{\lambda}$).
\item The objective $F(Q)=\sum_hQ(h)\RS(h)+\frac1{\lambda m}\sum_hQ(h)\ln\frac{Q(h)}{P(h)}$ is strictly convex
      ($x\mapsto x\ln x$ is strictly convex). With the Lagrangian $\mathcal L(Q,\nu)=F(Q)+\nu(\sum_hQ(h)-1)$,
      \[\frac{\partial\mathcal L}{\partial Q(h)}=\RS(h)+\frac{1}{\lambda m}\Big(\ln\frac{Q(h)}{P(h)}+1\Big)+\nu=0
      \iff Q(h)=P(h)\,e^{-\lambda m\RS(h)}\,e^{-1-\lambda m\nu}.\]
      The constant $e^{-1-\lambda m\nu}$ is fixed by $\sum_hQ(h)=1$, which gives the Gibbs posterior. The positivity
      constraints are automatically satisfied, and by strict convexity this stationary point is the unique minimum.
\end{enumerate}
\end{solution}

%% =====================================================================
\begin{exo}{From the general theorem to classical bounds}\label{exo:bounds}
Recall the general theorem: for convex $\Delta$, w.p.\ $\geq1-\delta$, for all $Q$,
$\Delta(\RS(\GQ),\RD(\GQ))\leq\frac1m[\KL(Q\|P)+\ln\frac{1}{\delta}\E_{S'}\E_{h\sim P}e^{m\Delta(\widehat R_{S'}(h),\RD(h))}]$.
\begin{enumerate}
\item Explain precisely where each of the following is used in the proof: convexity of $\Delta$, the change of measure,
      Markov's inequality, the independence of $P$ from $S$.
\item \emph{(Hoeffding-type bound).} Take $\Delta(q,p)=\lambda(p-q)$ for a fixed $\lambda>0$. Using Hoeffding's lemma, show that
      w.p.\ $\geq1-\delta$, for all $Q$: $\RD(\GQ)\leq\RS(\GQ)+\frac{\KL(Q\|P)+\ln\frac1\delta}{\lambda m}+\frac\lambda8$.
\item Minimize the right-hand side in $\lambda$. Why is it not allowed to plug this value into the bound? Propose a
      correct way to optimize $\lambda$ and give the resulting price.
\item \emph{(Catoni's bound).} Let $C>0$ and $\mathcal F_C(p)=-\ln(1-p(1-e^{-C}))$. Show that for a fixed voter with risk $p$ and
      the 0-1 loss, $\E_{S'}e^{-Cm\widehat R_{S'}(h)}=(1-p(1-e^{-C}))^m$. Deduce Catoni's bound.
\item Numerical application: $m=1000$, $\delta=0.05$, $\RS(\GQ)=0.1$, $\KL(Q\|P)=5$. Compute McAllester's bound
      $\RS+\sqrt{(\KL+\ln\frac{2\sqrt m}\delta)/(2m)}$ and the Hoeffding-type bound for $\lambda\in\{0.1,0.2\}$.
\end{enumerate}
\end{exo}

\begin{solution}
\begin{enumerate}
\item The proof has three steps (Section 3.2 of the notes).
      \emph{Convexity of $\Delta$} is used in the first step (Jensen): the risks of $\GQ$ are averages over $Q$ of the
      risks of the voters, and $\Delta(\E_Q\RS(h),\E_Q\RD(h))\leq\E_Q\Delta(\RS(h),\RD(h))$. This reduces the problem to
      single voters.
      \emph{The change of measure} is the second step: $\E_Q[m\Delta]\leq\KL(Q\|P)+\ln\E_Pe^{m\Delta}$. It removes $Q$
      from the random part of the inequality; since it holds for all $Q$ at once, so will the final bound.
      \emph{Markov's inequality} is the third and only probabilistic step: the non-negative variable
      $Z(S)=\E_Pe^{m\Delta(\RS(h),\RD(h))}$ satisfies $Z(S)\leq\E_{S'}Z(S')/\delta$ with probability $\geq1-\delta$.
      \emph{Independence of $P$ and $S$} is used inside this step, to exchange $\E_{S'}$ and $\E_{h\sim P}$ (Fubini):
      for each fixed $h$, $\widehat R_{S'}(h)$ is an average of i.i.d.\ variables, whose exponential moment can be
      computed. If $P$ depended on $S'$, this exchange would be meaningless.
\item Here $\Delta(q,p)=\lambda(p-q)$ is linear, hence convex. For a fixed $h$, $\widehat R_{S'}(h)=\frac1m\sum_i\ell_i$ with
      i.i.d.\ $\ell_i\in[0,1]$ of mean $\RD(h)$, so by independence
      \[\E_{S'}e^{m\lambda(\RD(h)-\widehat R_{S'}(h))}=\prod_{i=1}^m\E\,e^{\lambda(\RD(h)-\ell_i)}\leq\big(e^{\lambda^2/8}\big)^m,\]
      by Hoeffding's lemma (a centred variable with range $1$ has $\E e^{sX}\leq e^{s^2/8}$). Hence
      $\mathcal I_\Delta(m)\leq e^{m\lambda^2/8}$, and the theorem gives
      $\lambda(\RD(\GQ)-\RS(\GQ))\leq\frac1m\big[\KL(Q\|P)+\ln\frac1\delta\big]+\frac{\lambda^2}8$. Dividing by $\lambda$ gives the result.
\item Let $A=\frac{\KL(Q\|P)+\ln(1/\delta)}{m}$. The function $\lambda\mapsto\frac A\lambda+\frac\lambda8$ is minimal at
      $\lambda^\star=\sqrt{8A}$, where it equals $\sqrt{A/2}$: we recover a McAllester-type bound
      $\RS(\GQ)+\sqrt{\frac{\KL+\ln(1/\delta)}{2m}}$. \emph{Why it is forbidden:} $\lambda$ is a parameter of $\Delta$, i.e.\ of
      the probabilistic event on which the bound holds; this event must be fixed before drawing $S$. But
      $\lambda^\star$ depends on $\KL(Q\|P)$, hence on $Q$ and on $S$. Using it amounts to choosing, after seeing the
      data, among infinitely many events, each of probability $\geq1-\delta$: nothing guarantees that the chosen one
      holds. \emph{Correct way:} fix a grid $\lambda_1<\dots<\lambda_k$ in advance, apply the bound to each $\lambda_j$ with
      confidence $\delta/k$ (union bound): with probability $\geq1-\delta$ all $k$ bounds hold, and we can take the best
      one. The price is $\frac{\ln k}{\lambda m}$ (i.e.\ $\ln k$ added to $\ln\frac1\delta$) plus the error due to the
      discretization; a geometric grid with ratio $2$ covering $[10^{-3},1]$ has $k=10$ points and costs $\ln10\approx2.3$ nats.
\item For the 0-1 loss, $m\widehat R_{S'}(h)\sim\mathrm{Bin}(m,p)$ with $p=\RD(h)$, whose moment generating function is
      $\E e^{sB}=(1-p+pe^s)^m$. At $s=-C$: $\E_{S'}e^{-Cm\widehat R_{S'}(h)}=(1-p+pe^{-C})^m=(1-p(1-e^{-C}))^m=e^{-m\mathcal F_C(p)}$.
      Take $\Delta(q,p)=\mathcal F_C(p)-Cq$ (convex: $\mathcal F_C$ is convex, being $-\ln$ of an affine decreasing function,
      and $-Cq$ is linear). Then $\E_{S'}e^{m\Delta(\widehat R_{S'}(h),p)}=e^{m\mathcal F_C(p)}e^{-m\mathcal F_C(p)}=1$, so
      $\mathcal I_\Delta(m)=1$ and, with probability $\geq1-\delta$, for all $Q$:
      $\mathcal F_C(\RD(\GQ))\leq C\RS(\GQ)+\frac{\KL(Q\|P)+\ln(1/\delta)}m$. Since $\mathcal F_C$ is increasing with inverse
      $\mathcal F_C^{-1}(u)=\frac{1-e^{-u}}{1-e^{-C}}$, we obtain
      $\RD(\GQ)\leq\frac{1}{1-e^{-C}}\big[1-\exp\big(-C\RS(\GQ)-\frac{\KL+\ln(1/\delta)}m\big)\big]$. (For losses in $[0,1]$
      that are not binary, the equality becomes an inequality by convexity of $u\mapsto e^{-Cu}$.)
\item $\ln\frac{2\sqrt{1000}}{0.05}=\ln1264.9\approx7.14$, so McAllester's bound is
      $0.1+\sqrt{\frac{5+7.14}{2000}}=0.1+0.078=0.178$. For the Hoeffding-type bound, $\ln\frac1{0.05}\approx3.00$ and
      $\KL+\ln\frac1\delta=8.00$. With $\lambda=0.1$: $0.1+\frac{8.00}{100}+\frac{0.1}8=0.1+0.080+0.0125=0.1925$. With
      $\lambda=0.2$: $0.1+\frac{8.00}{200}+0.025=0.165$. The optimal value would be $\lambda^\star=\sqrt{8\times0.008}=0.253$,
      giving $0.1+\sqrt{0.004}=0.163$. Remember that $\lambda=0.2$ is legitimate only if it was chosen before seeing the
      data; if it was selected among $\{0.1,0.2\}$ after computing both bounds, each must be computed with $\delta/2$,
      which gives $0.1+\frac{5+\ln40}{200}+0.025=0.1+0.0434+0.025\approx0.168$ for $\lambda=0.2$.
\end{enumerate}
\end{solution}

%% =====================================================================
\begin{exo}{Occam's razor as a PAC-Bayesian bound}\label{exo:occam}
\begin{enumerate}
\item Let $\Hcal$ be countable and $Q=\delta_{h}$ a Dirac on $h\in\Hcal$. Compute $\KL(\delta_h\|P)$. What does Seeger's bound give
      for a single voter? Compare with the Occam bound of the lecture notes.
\item Why can we not use Dirac posteriors when $\Hcal$ is uncountable (e.g.\ linear classifiers with a Gaussian prior)?
\item Let $\Hcal=\{h_1,h_2,\dots\}$ be countable and $P(h_k)=\frac{6}{\pi^2k^2}$. Check that $P$ is a probability and give
      the complexity term of $h_k$. Interpret this prior in terms of \emph{description length}.
\end{enumerate}
\end{exo}

\begin{solution}
\begin{enumerate}
\item $\KL(\delta_h\|P)=\sum_{h'}\delta_h(h')\ln\frac{\delta_h(h')}{P(h')}=1\cdot\ln\frac1{P(h)}$. Seeger's bound with
      $Q=\delta_h$ (for which $\RS(G_{\delta_h})=\RS(h)$) gives, simultaneously for all $h$,
      \[\RD(h)\leq\klup\Big(\RS(h),\frac{\ln\frac1{P(h)}+\ln\frac{2\sqrt m}\delta}{m}\Big).\]
      The Occam bound of the notes has $\ln\frac1\delta$ instead of $\ln\frac{2\sqrt m}\delta$: it is slightly tighter
      because it only covers Dirac posteriors (a union bound over countably many events), whereas the PAC-Bayesian
      bound covers all posteriors simultaneously. For $m=1000$ the difference is $\ln(2\sqrt{1000})\approx4.1$ nats.
\item If $P$ has a density (e.g.\ a Gaussian on $\R^d$), then $P(\{h\})=0$ for every $h$, and $\delta_h$ is not
      absolutely continuous with respect to $P$: $\KL(\delta_h\|P)=+\infty$ and the bound is void. One must use
      ``spread'' posteriors, e.g.\ $\mathcal N(\mathbf w,\sigma^2I)$ around the chosen classifier; the bound then
      certifies the Gibbs classifier around $\mathbf w$, and the margin or the majority vote relates it to $\mathbf w$
      itself (Exercise~\ref{exo:linear}).
\item $\sum_{k\geq1}\frac1{k^2}=\frac{\pi^2}6$, so $\sum_kP(h_k)=1$ and $P$ is a probability. The complexity of $h_k$ is
      $\ln\frac1{P(h_k)}=2\ln k+\ln\frac{\pi^2}6\approx2\ln k+0.50$. \emph{Description length:} by Kraft's inequality, any
      prefix code assigning $\ell(h)$ nats to $h$ defines a sub-probability $P(h)=e^{-\ell(h)}$, and conversely. The
      prior thus encodes a way of describing hypotheses; the bound says that a hypothesis is cheap if it has a short
      description \emph{chosen before seeing the data}. Here the index $k$ is described with about $2\ln k$ nats, so
      hypotheses listed early are favoured: this is Occam's razor made quantitative.
\end{enumerate}
\end{solution}

%% =====================================================================
\begin{exo}{Bounds on the risk of the majority vote}\label{exo:mv}
We assume voters in $\{-1,+1\}$, 0-1 loss, and count ties as errors.
\begin{enumerate}
\item Show that $\RD(\BQ)=\P(W_Q\geq\frac12)$ and $\RD(\GQ)=\E W_Q$. Deduce $\RD(\BQ)\leq2\RD(\GQ)$. Give an
      example of distribution of $W_Q$ for which this inequality is an equality.
\item Show that $e_Q=\E[W_Q^2]$, $d_Q=\E[2W_Q(1-W_Q)]$, and $\RD(\GQ)=e_Q+\frac12d_Q$. Show also that
      $d_Q\leq2\RD(\GQ)(1-\RD(\GQ))$.
\item Prove the tandem bound $\RD(\BQ)\leq4e_Q$. Show that it is better than the first-order bound iff $d_Q\geq\RD(\GQ)$.
\item For $\mu<\frac12$, show that $\RD(\BQ)\leq\frac{\E(W_Q-\mu)^2}{(1/2-\mu)^2}$. Find the optimal $\mu^\star$ and show that
      the optimal value is $\frac{\Var W_Q}{\Var W_Q+(1/2-\E W_Q)^2}$ (Cantelli's inequality).
\item Deduce the C-bound $\RD(\BQ)\leq1-\frac{(1-2\RD(\GQ))^2}{1-2d_Q}$ (when $\RD(\GQ)<\frac12$).
\item \emph{(Condorcet).} $Q$ is uniform over $n$ voters erring independently with probability $p$. Compute $\RD(\GQ)$,
      $e_Q$, $d_Q$ as functions of $n$ and $p$. For $p=0.3$, compute the exact risk of the vote and the three bounds
      for $n=3$, then the C-bound for $n=25$ and its limit when $n\to\infty$. Comment.
\end{enumerate}
\end{exo}

\begin{solution}
\begin{enumerate}
\item For a fixed $(x,y)$, $y\,\E_Qh(x)=\E_Q[yh(x)]=\E_Q[1-2\ind{h(x)\neq y}]=1-2W_Q(x,y)$, because $yh(x)=1$ when
      $h$ is right and $-1$ when it errs. The vote errs (ties counted as errors) iff $y\,\E_Qh(x)\leq0$, i.e.\ iff
      $W_Q(x,y)\geq\frac12$. Taking the probability over $(x,y)\sim\D$: $\RD(\BQ)=\P(W_Q\geq\frac12)$. By Fubini,
      $\E_\D W_Q=\E_Q\E_\D\ind{h(x)\neq y}=\E_Q\RD(h)=\RD(\GQ)$. Markov's inequality applied to $W_Q\geq0$ gives
      $\P(W_Q\geq\frac12)\leq\frac{\E W_Q}{1/2}=2\RD(\GQ)$.
      \emph{Equality case:} Markov is tight when $W_Q$ only takes the values $0$ and $\frac12$. Example: two voters
      with weight $\frac12$ each, which never err at the same time, and such that on a fraction $r$ of the points
      exactly one of them errs. Then $W_Q=\frac12$ on these points, the vote errs on them (tie), $\RD(\BQ)=r$ and
      $\RD(\GQ)=\frac r2$.
\item For fixed $(x,y)$, draw $h,h'\sim Q$ independently. The events ``$h$ errs'' and ``$h'$ errs'' are independent
      with probability $W_Q$ each, so both err with probability $W_Q^2$, and exactly one errs with probability
      $2W_Q(1-W_Q)$; for binary voters, $h$ and $h'$ disagree iff exactly one of them errs. Averaging over $\D$ (Fubini):
      $e_Q=\E W_Q^2$ and $d_Q=\E[2W_Q(1-W_Q)]$. Pointwise $W_Q=W_Q^2+\frac12\cdot2W_Q(1-W_Q)$, hence
      $\RD(\GQ)=e_Q+\frac12d_Q$. Finally Jensen (or $\Var W_Q\geq0$) gives $e_Q=\E W_Q^2\geq(\E W_Q)^2=\RD(\GQ)^2$,
      so $d_Q=2(\RD(\GQ)-e_Q)\leq2\RD(\GQ)(1-\RD(\GQ))$, with equality iff $W_Q$ is almost surely constant.
\item $\P(W_Q\geq\frac12)=\P(W_Q^2\geq\frac14)\leq\frac{\E W_Q^2}{1/4}=4e_Q$ (Markov applied to $W_Q^2$). Using question~2,
      $4e_Q=4\RD(\GQ)-2d_Q$, so $4e_Q\leq2\RD(\GQ)\iff2\RD(\GQ)\leq2d_Q\iff d_Q\geq\RD(\GQ)$: the tandem bound improves on
      the first-order bound exactly when the voters disagree more often than an average voter errs.
\item For $\mu<\frac12$ and $W_Q\geq\frac12$, we have $W_Q-\mu\geq\frac12-\mu>0$, hence $(W_Q-\mu)^2\geq(\frac12-\mu)^2$: the
      event $\{W_Q\geq\frac12\}$ is included in $\{(W_Q-\mu)^2\geq(\frac12-\mu)^2\}$ and Markov gives the bound. Write $R=\E W$,
      $e=\E W^2$, $v=e-R^2$ and $f(\mu)=\frac{e-2\mu R+\mu^2}{(1/2-\mu)^2}$. Then
      \[f'(\mu)=\frac{2(\mu-R)(\frac12-\mu)+2(e-2\mu R+\mu^2)}{(1/2-\mu)^3}=\frac{2\big(e-\frac R2-\mu(\frac12-R)\big)}{(1/2-\mu)^3},\]
      which vanishes at $\mu^\star=\frac{R/2-e}{1/2-R}$ (when $R<\frac12$), is negative before and positive after: $\mu^\star$ is
      the minimum. To compute $f(\mu^\star)$, note that $\mu^\star-R=\frac{R/2-e-R/2+R^2}{1/2-R}=\frac{-v}{1/2-R}$ and
      $\frac12-\mu^\star=\frac{1/4-R/2-R/2+e}{1/2-R}=\frac{v+(1/2-R)^2}{1/2-R}$. The numerator is
      $e-2\mu^\star R+\mu^{\star2}=(\mu^\star-R)^2+v=\frac{v^2}{(1/2-R)^2}+v=\frac{v\,(v+(1/2-R)^2)}{(1/2-R)^2}$, so
      \[f(\mu^\star)=\frac{v\,(v+(1/2-R)^2)}{(1/2-R)^2}\cdot\frac{(1/2-R)^2}{(v+(1/2-R)^2)^2}=\frac{v}{v+(1/2-R)^2},\]
      which is Cantelli's inequality for $W_Q$ with $a=\frac12-R$.
\item With $M=1-2W_Q$ (the margin), $\Var M=4v$ and $\E M=1-2R$, so $\frac{v}{v+(1/2-R)^2}=\frac{4v}{4v+(1-2R)^2}=\frac{\Var M}{\E M^2}=1-\frac{(\E M)^2}{\E M^2}$.
      Finally $\E M=1-2\RD(\GQ)$ and $\E M^2=1-4\E W+4\E W^2=1-4\RD(\GQ)+4e_Q=1-2d_Q$ (question~2). Hence
      $\RD(\BQ)\leq1-\frac{(1-2\RD(\GQ))^2}{1-2d_Q}$: the C-bound is the best member of the family of question~4.
\item $nW_Q\sim\mathrm{Bin}(n,p)$ since each voter errs independently with probability $p$. Hence $\RD(\GQ)=\E W=p$,
      $e_Q=\E W^2=\Var W+p^2=\frac{p(1-p)}n+p^2$, and $d_Q=2(p-e_Q)=2p(1-p)(1-\frac1n)$.
      For $p=0.3$ and $n=3$, the vote errs iff at least $2$ voters err: $\RD(\BQ)=3p^2(1-p)+p^3=0.189+0.027=0.216$. The
      bounds are: first order $2p=0.6$; $e_Q=0.09+0.07=0.16$, so tandem $=0.64$; $d_Q=2\times0.21\times\frac23=0.28$, so
      C-bound $=1-\frac{(0.4)^2}{1-0.56}=1-\frac{0.16}{0.44}\approx0.636$. For $n=25$: $d_Q=0.42\times0.96=0.4032$,
      $1-2d_Q=0.1936$ and C-bound $=1-\frac{0.16}{0.1936}\approx0.174$, while the true risk is about $0.018$. When $n\to\infty$,
      $1-2d_Q\to1-4p(1-p)=(1-2p)^2$ and the C-bound tends to $0$, like the true risk, at the rate $\frac{4p(1-p)}{n(1-2p)^2}$.
      \emph{Comment:} with $3$ voters, $d_Q=0.28<p=0.3$, and the second-order bounds are worse than the first-order
      one (Proposition 5.7 of the notes); with many independent voters, only the second-order bounds, and
      especially the C-bound, capture the benefit of the vote. The tandem bound tends to $4p^2=0.36$ and never goes
      to $0$.
\end{enumerate}
\end{solution}

%% =====================================================================
\begin{exo}{Linear classifiers with a Gaussian posterior}\label{exo:linear}
Voters $h_{\mathbf v}(x)=\sign\langle\mathbf v,x\rangle$, prior $P=\mathcal N(\mathbf 0,I_d)$, posterior $Q_{\mathbf w}=\mathcal N(\mathbf w,I_d)$.
\begin{enumerate}
\item Show that $\BQ(x)=\sign\langle\mathbf w,x\rangle$.
\item Show that $\P_{\mathbf v\sim Q_{\mathbf w}}\big(h_{\mathbf v}(x)\neq y\big)=\Phi\big(-\frac{y\langle\mathbf w,x\rangle}{\|x\|}\big)$ where $\Phi$ is the standard
      Gaussian CDF. Give $\RS(G_{Q_{\mathbf w}})$ and $\KL(Q_{\mathbf w}\|P)$.
\item Show that, for a fixed $C$, minimizing Catoni's bound over $\mathbf w$ amounts to minimizing
      $C\sum_i\Phi(-y_i\langle\mathbf w,x_i\rangle/\|x_i\|)+\frac12\|\mathbf w\|^2$. Compare with the SVM.
\item For $\mathbf w=s\mathbf u$ with $\|\mathbf u\|=1$, describe the evolution of $\RS(G_{Q_{s\mathbf u}})$ and of the KL when $s$ grows.
      Relate to the notion of margin.
\end{enumerate}
\end{exo}

\begin{solution}
\begin{enumerate}
\item Fix $x$. Under $Q_{\mathbf w}$, $\langle\mathbf v,x\rangle$ is Gaussian with mean $\langle\mathbf w,x\rangle$ and variance
      $x^\top I_dx=\|x\|^2$. Hence $\P(\langle\mathbf v,x\rangle>0)=\Phi\big(\frac{\langle\mathbf w,x\rangle}{\|x\|}\big)$ and
      $\E_Q\sign\langle\mathbf v,x\rangle=2\Phi\big(\frac{\langle\mathbf w,x\rangle}{\|x\|}\big)-1$, which is positive iff
      $\langle\mathbf w,x\rangle>0$ (as $\Phi(0)=\frac12$ and $\Phi$ is increasing). Therefore $\BQ(x)=\sign\langle\mathbf w,x\rangle$:
      the vote of infinitely many noisy linear classifiers is the linear classifier at their centre.
\item Similarly $y\langle\mathbf v,x\rangle\sim\mathcal N(y\langle\mathbf w,x\rangle,\|x\|^2)$ (as $y^2=1$), so
      $\P_{\mathbf v}(h_{\mathbf v}(x)\neq y)=\P(y\langle\mathbf v,x\rangle\leq0)=\Phi\big(-\frac{y\langle\mathbf w,x\rangle}{\|x\|}\big)$.
      Averaging over the sample, $\RS(G_{Q_{\mathbf w}})=\frac1m\sum_i\Phi\big(-\frac{y_i\langle\mathbf w,x_i\rangle}{\|x_i\|}\big)$: a
      smooth \emph{probit} loss of the normalized margin $\frac{y_i\langle\mathbf w,x_i\rangle}{\|x_i\|}$. By
      Exercise~\ref{exo:kl}.4 with $\sigma=1$, $\KL(Q_{\mathbf w}\|P)=\frac12\|\mathbf w\|^2$.
\item Catoni's bound reads $\RD(\GQ)\leq\mathcal F_C^{-1}\big(C\RS(\GQ)+\frac{\KL(Q\|P)+\ln(1/\delta)}m\big)$, and
      $\mathcal F_C^{-1}$ is increasing: for fixed $C$ and $\delta$, minimizing the bound over $\mathbf w$ amounts to minimizing
      $C\RS(G_{Q_{\mathbf w}})+\frac{\|\mathbf w\|^2}{2m}$, i.e.\ (multiplying by $m$)
      $C\sum_i\Phi\big(-\frac{y_i\langle\mathbf w,x_i\rangle}{\|x_i\|}\big)+\frac12\|\mathbf w\|^2$. The soft-margin SVM minimizes
      $C\sum_i\max(0,1-y_i\langle\mathbf w,x_i\rangle)+\frac12\|\mathbf w\|^2$: the regularizer is the same (it is the KL to the
      prior), and the hinge loss is replaced by the probit loss. Two differences: the probit loss is bounded (a badly
      misclassified point costs at most $1$, which gives robustness to label noise) but not convex; and it uses the
      normalized margin.
\item For $\mathbf w=s\mathbf u$, the loss of point $i$ is $\Phi(-s\gamma_i)$ with $\gamma_i=\frac{y_i\langle\mathbf u,x_i\rangle}{\|x_i\|}$.
      When $s$ grows, it decreases to $0$ for well-classified points ($\gamma_i>0$), the faster the larger $\gamma_i$, and
      increases to $1$ for misclassified points; at $s=0$ all losses equal $\frac12$. Meanwhile the KL grows as $\frac{s^2}2$.
      If all points have a normalized margin at least $\gamma$, choosing $s\approx\frac{c}{\gamma}$ makes the empirical Gibbs
      risk at most $\Phi(-c)$ at a KL cost of $\frac{c^2}{2\gamma^2}$: large margins allow small KL, which is the
      PAC-Bayesian margin bound (Langford and Shawe-Taylor, 2002). The trade-off in $s$ is the one of Figure~31 of the notes (PBGD).
\end{enumerate}
\end{solution}

%% =====================================================================
\begin{exo}{$(\star)$ The PAC-Bayes-$\lambda$ bound and its minimization}\label{exo:lambda}
\begin{enumerate}
\item Prove the refined Pinsker inequality $\kl(q\|p)\geq\frac{(p-q)^2}{2p}$ for $q<p$.
\item Starting from Seeger's bound, show that w.p.\ $\geq1-\delta$, for all $Q$ and all $\lambda\in(0,2)$ simultaneously,
      $\RD(\GQ)\leq\frac{\RS(\GQ)}{1-\lambda/2}+\frac{\KL(Q\|P)+\ln\frac{2\sqrt m}\delta}{\lambda(1-\lambda/2)m}$.
      Why is no union bound over $\lambda$ needed here, contrary to Exercise~\ref{exo:bounds}?
\item For a fixed $Q$, compute the optimal $\lambda$. For a fixed $\lambda$, compute the optimal $Q$ on a finite set of voters.
      Write the alternating minimization algorithm.
\item Explain why the posterior returned by this algorithm on a random forest may give a \emph{worse} majority vote than the uniform one.
\end{enumerate}
\end{exo}

\begin{solution}
\begin{enumerate}
\item Let $0\leq q<p\leq1$ and $g(u)=\kl(u\|p)-\frac{(p-u)^2}{2p}$ for $u\in[0,p]$. Then $g(p)=0$ and
      $g'(u)=\ln\frac{u(1-p)}{p(1-u)}+\frac{p-u}{p}$. Using $\ln t\leq t-1$ with $t=\frac{u(1-p)}{p(1-u)}$:
      $\ln\frac{u(1-p)}{p(1-u)}\leq\frac{u(1-p)-p(1-u)}{p(1-u)}=\frac{u-p}{p(1-u)}\leq\frac{u-p}{p}$, the last step because
      $u-p<0$ and $0<1-u\leq1$. Hence $g'(u)\leq0$: $g$ is non-increasing on $(0,p)$ and $g(q)\geq g(p)=0$ (see also
      Appendix A.3 of the notes).
\item Let $\hat q=\RS(\GQ)$, $p=\RD(\GQ)$ and $\varepsilon=\frac1m[\KL(Q\|P)+\ln\frac{2\sqrt m}\delta]$. On the event of Seeger's bound,
      $\kl(\hat q\|p)\leq\varepsilon$. If $p\leq\hat q$ the claimed bound is trivial. Otherwise the refined Pinsker inequality
      gives $\frac{(p-\hat q)^2}{2p}\leq\varepsilon$, i.e.\ $p-\hat q\leq\sqrt{2p\varepsilon}$. For any $\lambda>0$, the
      arithmetic--geometric inequality $\sqrt{ab}\leq\frac{a+b}2$ with $a=\lambda p$, $b=\frac{2\varepsilon}\lambda$ gives
      $\sqrt{2p\varepsilon}\leq\frac{\lambda p}2+\frac\varepsilon\lambda$. Hence $p(1-\frac\lambda2)\leq\hat q+\frac\varepsilon\lambda$, and for
      $\lambda\in(0,2)$ we can divide by $1-\frac\lambda2>0$. \emph{Why no union bound:} all these inequalities are
      deterministic consequences of the \emph{single} event of Seeger's bound; $\lambda$ appears only in the algebra,
      not in the probabilistic statement. In Exercise~\ref{exo:bounds}, on the contrary, $\lambda$ defined the function
      $\Delta$, hence the event itself.
\item \emph{Optimal $\lambda$ for fixed $Q$.} Let $B(\lambda)=\frac{\hat q}{1-\lambda/2}+\frac{\varepsilon}{\lambda(1-\lambda/2)}$. Setting
      $B'(\lambda)=0$ leads to $\frac{\hat q}{2}\lambda^2=\varepsilon(1-\lambda)$, i.e.\ $\hat q\lambda^2+2\varepsilon\lambda-2\varepsilon=0$, whose
      positive root is $\lambda^\star=\frac{-\varepsilon+\sqrt{\varepsilon^2+2\hat q\varepsilon}}{\hat q}=\frac{2}{\sqrt{2\hat q/\varepsilon+1}+1}$.
      \emph{Optimal $Q$ for fixed $\lambda$.} The bound is an increasing affine function of
      $\RS(\GQ)+\frac{\KL(Q\|P)}{\lambda m}$, minimized by the Gibbs posterior $Q(h)\propto P(h)e^{-\lambda m\RS(h)}$
      (Exercise~\ref{exo:dv}). \emph{Algorithm:} $\lambda\leftarrow1$; repeat $Q\leftarrow$ Gibbs posterior at $\lambda$, then
      $\lambda\leftarrow\lambda^\star(Q)$, until $\lambda$ stabilizes; return $Q$ and the certificate computed with the kl form. Each
      step decreases the bound, which is therefore non-increasing along the iterations (and converges to the global
      minimum under the conditions of Thiemann et al., 2017).
\item The bound only involves the \emph{average} OOB loss $\rho^\top\widehat L$ and the KL. It therefore puts the weight
      on the few trees with the smallest OOB loss, which tend to be similar (they are grown on similar bootstrap
      samples and make similar errors). The effective number of trees collapses (to about $2$ on \texttt{digits}),
      the disagreement vanishes, and the vote loses the benefit of averaging: its risk approaches that of a single
      good tree ($0.150$ instead of $0.044$ on \texttt{digits}, Table 2 of the notes). The bound is excellent for the
      Gibbs classifier, but it is the wrong objective for the vote.
\end{enumerate}
\end{solution}

%% =====================================================================
\begin{exo}{$(\star)$ Out-of-bag samples}\label{exo:oob}
Consider bagging with bootstrap samples of size $k$ drawn with replacement among the $m$ training examples.
\begin{enumerate}
\item Compute the probability that a given example is out-of-bag for a given tree; give its limit when $m\to\infty$ for
      $k=m$ and for $k=m/2$. Same question for being out-of-bag for two given trees.
\item Explain why the OOB loss $\widehat L_t$ of tree $t$ is an unbiased estimate of $\RD(h_t)$ given $S_t$, and why a uniform
      prior over the \emph{indices} of the trees is a legitimate PAC-Bayesian prior although the trees depend on $S$.
\item Show that if $\E e^{N\kl(\hat p_N\|p)}\leq2\sqrt N$ for all $N$, then for $n_0\leq N$, $\E e^{n_0\kl(\hat p_N\|p)}\leq2\sqrt{n_0}$.
      Why is this useful?
\end{enumerate}
\end{exo}

\begin{solution}
\begin{enumerate}
\item A bootstrap sample of size $k$ draws $k$ indices uniformly and independently in $\{1,\dots,m\}$. A given example
      is never drawn with probability $(1-\frac1m)^k$. Since $(1-\frac1m)^m\to e^{-1}$, this tends to
      $e^{-1}\approx0.368$ for $k=m$ and to $e^{-1/2}\approx0.607$ for $k=m/2$. For two trees with independent bootstrap
      samples, the example is out-of-bag for both with probability $(1-\frac1m)^{2k}\to e^{-2}\approx0.135$ ($k=m$) or
      $e^{-1}\approx0.368$ ($k=m/2$). Halving the bootstrap size thus almost triples the size of the pairwise OOB sets
      used by the tandem bound, at the price of slightly weaker trees.
\item Condition on the bootstrap sample $S_t$ (and on the internal randomness of the tree). Then $h_t$ is a fixed
      function, and the OOB examples, which were not used to build it, are still i.i.d.\ from $\D$ and independent of
      $h_t$ (they are the examples of $S$ whose index was not drawn; given the drawn indices, their values are
      independent of $S_t$). Hence $\E[\widehat L_t\,|\,S_t]=\RD(h_t)$, and Maurer's lemma applies to $\widehat L_t$
      conditionally on $S_t$. The prior is placed on the \emph{indices} $t\in\{1,\dots,n\}$, not on the trees themselves:
      the uniform distribution on indices is fixed before seeing the data. The dependence of $h_t$ on $S$ only enters the
      third step of the proof, where, for each index $t$, the conditional concentration of $\widehat L_t$ is used.
\item Write $e^{n_0\kl(\hat p_N\|p)}=\big(e^{N\kl(\hat p_N\|p)}\big)^{n_0/N}$. Since $n_0\leq N$, $x\mapsto x^{n_0/N}$ is concave,
      and Jensen gives $\E e^{n_0\kl}\leq\big(\E e^{N\kl}\big)^{n_0/N}\leq(2\sqrt N)^{n_0/N}=\exp\big(n_0\frac{\ln(2\sqrt N)}{N}\big)$.
      The function $N\mapsto\frac{\ln(2\sqrt N)}N$ is decreasing for $N\geq1$ (its derivative is $\frac{1/2-\ln(2\sqrt N)}{N^2}<0$),
      so for $N\geq n_0$ this is at most $\exp\big(n_0\frac{\ln(2\sqrt{n_0})}{n_0}\big)=2\sqrt{n_0}$. \emph{Usefulness:} the OOB sets
      have different sizes; this lemma allows us to use a single sample size $n_0=\min_t|T_t|$ in the complexity term
      while still computing each loss $\widehat L_t$ on \emph{all} the OOB examples of tree $t$.
\end{enumerate}
\end{solution}


%% =====================================================================
\begin{exo}{$(\star)$ Self-bounding learning}\label{exo:selfbounding}
Let $\mathcal B(Q,S,\delta)$ be a bound such that $\P_S(\forall Q,\ \RD(Q)\leq\mathcal B(Q,S,\delta))\geq1-\delta$.
\begin{enumerate}
\item Let $Q_S$ be the posterior returned by any algorithm run on $S$. Show that $\P_S(\RD(Q_S)\leq\mathcal B(Q_S,S,\delta))\geq1-\delta$.
      Why does this fail for a bound that holds only for a \emph{fixed} predictor (e.g.\ Hoeffding's inequality)?
\item You build $K=5$ random forests with tree depths $1,2,4,8$ and unlimited, and compute for each of them the OOB tandem
      certificate. How must you modify the certificates to be allowed to keep the forest with the smallest one?
      What is the cost in the complexity term?
\item Among the following choices, which are free, which require a union bound, which are forbidden?
      (a) the temperature $\lambda$ of a Gibbs posterior; (b) the parameter $C$ of Catoni's bound; (c) the variance of a
      Gaussian prior chosen on a grid of 10 values; (d) the centre of a Gaussian prior learned on the sample used in the
      bound; (e) the number of gradient iterations when minimizing the bound.
\item Let $p=\klup(q,\varepsilon)$. Show that $\frac{\partial p}{\partial\varepsilon}=\frac{p(1-p)}{p-q}$ and
      $\frac{\partial p}{\partial q}=\frac{p(1-p)}{p-q}\ln\frac{p(1-q)}{q(1-p)}$. Deduce the gradient of the first-order
      certificate $2\,\klup\big(\rho^\top\widehat L,\frac{\KL(\rho\|\pi)+c}{n}\big)$ with respect to $\rho$.
\item Why is the test-set bound usually tighter than a self-bounding certificate for a random forest? Give one
      advantage of the self-bounding approach.
\end{enumerate}
\end{exo}

\begin{solution}
\begin{enumerate}
\item Let $E=\{S:\forall Q,\ \RD(Q)\leq\mathcal B(Q,S,\delta)\}$, with $\P(E)\geq1-\delta$. For $S\in E$ the inequality holds for
      every posterior, in particular for $Q_S$, whatever the way it was computed. Hence
      $\{S:\RD(Q_S)\leq\mathcal B(Q_S,S,\delta)\}\supseteq E$ and has probability $\geq1-\delta$. For a fixed-predictor bound,
      there is one event $E_h$ per predictor, each of probability $\geq1-\delta$, but their intersection can have a much
      smaller probability. The learning algorithm, which looks at $S$, may precisely choose a predictor for which
      $E_h$ fails: among $1000$ random classifiers of risk $0.5$ evaluated on $100$ examples, the one with the smallest
      empirical risk violates the naive Hoeffding bound in $99.8\%$ of the runs (Figure 10 of the notes).
\item Keeping the smallest certificate is a data-dependent choice among $K=5$ bounds. By the union bound
      (Proposition 6.3 of the notes), if each certificate is computed with confidence $\delta/5$, all five hold
      simultaneously with probability $\geq1-\delta$, and so does the smallest one. For the tandem bound, the complexity
      term $\frac{2\KL+\ln(2\sqrt{n_0}/\delta)}{n_0}$ becomes $\frac{2\KL+\ln(2\sqrt{n_0}/\delta)+\ln5}{n_0}$: the price is
      $\ln5\approx1.6$ nats, i.e.\ $0.004$ in the complexity term for $n_0=400$.
\item (a) \emph{Free}: a temperature only indexes a family of posteriors compared with the same prior in the same
      bound. (b) \emph{Union bound}: $C$ defines the function $\Delta$, hence the event; use a grid fixed in advance with
      $\delta/k$. (c) \emph{Union bound}: each variance defines a different prior, hence a different bound; cost $\ln10$.
      (d) \emph{Forbidden}: the prior would depend on the sample used in the bound, and step 3 of the proof fails; it
      becomes legitimate if the centre is learned on a separate part $S_1$ and the bound computed on $S_2$
      (Proposition 6.5). (e) \emph{Free}: any stopping rule produces a posterior, and the bound holds for all posteriors.
\item Let $F(q,p)=\kl(q\|p)-\varepsilon$, so that $F(q,\klup(q,\varepsilon))=0$. We have $F_q=\partial_q\kl(q\|p)=\ln\frac{q(1-p)}{p(1-q)}$
      and $F_p=\frac{p-q}{p(1-p)}\neq0$ (since $p>q$ when $\varepsilon>0$). By the implicit function theorem,
      $\frac{\partial p}{\partial\varepsilon}=-\frac{F_\varepsilon}{F_p}=\frac1{F_p}=\frac{p(1-p)}{p-q}$ and
      $\frac{\partial p}{\partial q}=-\frac{F_q}{F_p}=\frac{p(1-p)}{p-q}\ln\frac{p(1-q)}{q(1-p)}$. For the first-order certificate
      $B(\rho)=2\,\klup(q(\rho),\varepsilon(\rho))$ with $q(\rho)=\rho^\top\widehat L$ and $\varepsilon(\rho)=\frac{\KL(\rho\|\pi)+c}n$,
      we have $\nabla_\rho q=\widehat L$ and $\nabla_\rho\varepsilon=\frac1n\big(\ln\frac\rho\pi+\mathbf1\big)$, so by the chain rule
      \[\nabla_\rho B=2\Big[\frac{\partial p}{\partial q}\,\widehat L+\frac{\partial p}{\partial\varepsilon}\,\frac1n\Big(\ln\frac\rho\pi+\mathbf 1\Big)\Big].\]
      (With a softmax parametrization $\rho=\mathrm{softmax}(\theta)$, multiply by the Jacobian: $\nabla_\theta B=\rho\odot(g-\rho^\top g)$.)
      Numerically, at $q=0.1$, $\varepsilon=0.05$, $p\approx0.220$: $\partial_\varepsilon p\approx1.43$ and $\partial_qp\approx1.33$.
\item The test-set bound certifies the vote itself, with a single kl inversion, no KL term, and no factor $2$ or $4$
      coming from the Gibbs-to-vote relation; on a random forest this relation is the main source of looseness. In
      exchange, the test examples are not used for learning, and a new split is needed for model selection. The
      self-bounding approach uses all the data, provides a learning objective with a guarantee, and allows model
      selection at the price of $\ln K$; with data-dependent priors it can be competitive when enough data are
      available.
\end{enumerate}
\end{solution}

\end{document}
